Page 11, pdfpage 15
\begin{itemize}
\item[a)]
A group is
\underline{Abelian} if all its elements commute. It follows from the correspondence
between commutators and square brackets that for an Abelian group all
square brackets, and consequently all structure constants, vanish:
\begin{equation}\label{eq:20}
c^{\tau}_{\rho \sigma} =0
\end{equation}
\item[b)]
A
\underline{subgroup} of a group \(S\) is a subset of elements of \(S\) which
satisfies the group postulates. Thus, if \(X_1, X_2, \ldots, X_p\) are the
infinitesimal operators of a subgroup, the structure constant's of the
group must satisfy the relations
\begin{equation}\label{eq:21}
c_{\rho \sigma}^\tau=0 \quad \rho, \sigma \leq p, \quad p < \tau.
\end{equation}
\item[c)]
An \underline{invariant} \underline{subgroup}, \(H\), of a group \(S\) is
a subgroup of \(S\) which contains all the
conjugates (images) of its elements. Thus, with \(S_n\), it contains \(S_x
S_n S_x^{-1}\) for any \(S_x\) in \(S\). If so, it, also contains the
commutator \(\mathrm{S}_{\mathrm{x}} \mathrm{S}_{\mathrm{n}}
\mathrm{S}_{\mathrm{x}}^{-1} \mathrm{~S}_{\mathrm{n}}^{-1}\). Thus, the
square bracket connecting an infinitesimal element of \(H\) with any
infinitesimal element of \(S\) must belong to \(H\). If \(X_1, X_2,
\ldots, X_p\) are the infinitesimal operators of an invariant subgroup of
\(S\), the structure constants of \(S\) must satisfy
\begin{equation}\label{eq:22}
c_{\rho \sigma}^\tau=0 \quad \rho \leq p, \quad p < \tau.
\end{equation}
\item[d)]
A group is \underline{simple}
if it has no invariant subgroups besides the unit element.
\item[e)] A group is \underline{semi-simple}
if it has no Abelian invariant subgroups besides the
unit element.
\end{itemize}
Page 12 pdfpagae 16

The distinction between groups which have Abelian invariant subgroups and
those which do not have such subgroups is important, because Abelian
subgroups, though apparently easiest to deal with, can actually be most
troublesome from the point of view of representations, as the following
example will show:

We consider the group of rectilinear motions in one dimension, in which the transformation \(x!=x+a\) followed by \(x^{\prime \prime}=x^{\prime}+b\) is equivalent to \(x^{\prime \prime}=x+a+b\). This group can be represented by square matrices of the second rank, in terms of which the composition law just given would read
\[
\begin{bmatrix}
1 & a \\
0 & 1
\end{bmatrix}
\begin{bmatrix}
1 & b \\
0 & 1
\end{bmatrix}=
\begin{bmatrix}
1 & a+b \\
0 & 1
\end{bmatrix}
\]
However, none of the matrices of this particular representation can be
brought to diagonal form by a similarity transformation. This peculiar
behavior is closely related to the Abelian property. Indeed, as we shall
show later, semi-simple groups never exhibit it. Moreover, the physical
applications in which we shall be interested will require the use only of
semi-simple groups. We shall therefore from this point on restrict
ourselves to the study of semi-simple groups. To this end, we must have a
criterion for their identification.

Such a criterion can be formulated very simply in terms of a symmetrical
tensor of the second rank which we construct from the
\(c_{\rho \sigma}^\tau\):
\begin{equation}\label{eq:23}
g_{\rho \sigma}=c_\rho \lambda^\mu c_{\sigma \cdot \mu} \lambda .
\end{equation}

Page 13, pdf page 17

If the group is sami-simple, then necessarily,
\begin{equation}\label{eq:24}
\det g_{\rho \tau} \ne  0
\end{equation}
For suppose it possesses an Abelian invariant subgroup, the indices of whose
elements are denoted by \(\bar{\rho}, \bar{\sigma}, \ldots\). Then,
\[
\begin{aligned}
g_\rho & =c_\rho \mu^\mu c_{\bar{\sigma}} \lambda & & \\
& =c_\rho{ }^\mu c_\sigma \bar{\lambda} & & \text { by (22) } \\
& =c_\rho \bar{\mu}^{\prime} c_{\bar{\sigma} \bar{\mu}} \bar{\lambda} & & \text { by (22) } \\
& =0, & & \text { by (20), }
\end{aligned}
\]
That the condition (\ref{eq:24}) is sufficient as well as necessary has been shown by Cartan
%\cite{Cartan1894}
.


We can use the tensor \(g_{\mu \nu}\) to define a relation of orthogonality
between contravariant vectors or to form new tensors by lowering of indices. As
an example,
\begin{equation}\label{eq:25}
c_{\rho \sigma \lambda}=c_{\rho \sigma}^{\tau} g_{\tau \lambda}
\end{equation}
and this new tensor is totally antisymmetric, for by (23),
\[
\begin{array}{rlr}
c_{\rho \sigma \lambda} & =c_{\rho \sigma} \tau_\sigma{ }_{\tau \dot{\mu}}^\nu c_{\lambda \nu}^\mu, & \\
& =-c_{\sigma \mu} \tau_\tau c^\nu c_{\lambda \nu}^\mu-c_\mu c_{\tau \sigma} c_\tau c_{\lambda \nu}^\mu & \text { by }\left(c_2\right) \\
& =c_\sigma \tau_\sigma c_{\rho \tau}^\nu c_{\lambda \nu}{ }^\mu+{ }_{\mu \rho} c_\mu c_{\tau \sigma}{ }^\nu c_{\nu \lambda}^\mu & \text { by }\left(c_1\right) .
\end{array}
\]
The last line has the desiréd property, since it is invariant under cyclic permutation of the indices and is, by construction, skew in \(\rho\) and \(\sigma\).

Page 14

If the group \(S\) is semi-simple, then Cartan's criterion (24) implies that we
can form from \(g_{\rho \sigma}\) the reciprocal tensor \(g^{\rho \sigma}\)
which can be used to raise indices and define orthogonality between covariant vectors.

\textcolor{blue}{
I can not longer ignore the fact that I am lost.
Guess.  \(L_i\) is a rotation about the \(i\)th axis. 
\[
L_3= \begin{bmatrix}
0 & -1 & 0 \\
1 & 0  & 0\\
0 & 0  & 1
\end{bmatrix},
P = \begin{bmatrix}
0 & 1 & 0 \\
0 & 0 & 1\\
1 & 0 & 0
\end{bmatrix},
\]
\(L_{i+1}=P^T L_i P\).  I am making mistakes everywhere, but
I do not think that fixing them will salvage this approach.
\[
[L_1 L_2] = \begin{bmatrix}
o & -1 & 1 \\
1 & o  & 1 \\
1 & 1  & o 
\end{bmatrix}
[L_2 L_3] = \begin{bmatrix}
o & 1 & 1 \\
1 & o  & -1 \\
1 & 1  & o 
\end{bmatrix}
[L_3 L_1] = \begin{bmatrix}
o & 1 & -1 \\
1 & o  & 1 \\
1 & 1  & o 
\end{bmatrix}.
\]
This is not right in more ways that I can imagine.  \(k=i+2\),
\[
[L_i L_{i+1}] = L_{k} + e_k a_k' + a_k e_k',  a_k= (1,1,1) - e_k 3/2
\]
Oops.
\[
[L_3 L_1] = -L_2 +c e_2' + e_2 c', c=(1/2,1,1)
\]
What if \(L_i\) is a skew matrix?
cross(x,u) = S(x) u, \(S(x) = L_i x^{i}\)
\[
S(x) = 
\begin{bmatrix}
 o &-z & y\\
 z & o &-x\\
-y & x & o
\end{bmatrix}
\]
\(S(u)=L_j u^{j}\), \(S(x)S(u) = u x' - <x,u> I\)
\( L_i L_j x^i  u^j\)
}

As an example of the foregoing we consider the group of rigid motions in three
dimensions, consisting of rotations and translations. The infinitesimal
rotations are generated by operators \(L_j(j=1,2,3)\) satisfying
\[
\left[L_1 L_2\right]=i L_3 \text { etc }
\]
and the infinitesimal displacements by operators \(P_1=L_4, P_2=L_5, P_3=L_6\) which commute among themselves but which satisfy
\[
\left[L_1 L_5\right]=i L_6 \text { etc. }
\]
so that the only non-vanishing structure constents are
\[
c_{12}{ }^3=c_{23}{ }^1=c_{15}{ }^6=c_{26}{ }^4=c_{34}{ }^5=c_{31}{ }^2=c_{61}{ }^5=c_{42}{ }^6=c_{53}{ }^4=1
\]
plus a corresponding list given by \(\left(C_1\right)\). For the \(g_{\rho \sigma}\) we find
\[
g_{11}=g_{22}=g_{33}=4, g_{44}=g_{55}=g_{66}=0, g_{\rho \sigma}=0 \quad(\rho \neq \sigma) .
\]

The determinant det \(g_{\rho \sigma}\) vanishes, as required by Cartan's
criterion, since the translations form an Abelian invariant subgroup,

If we consider only the group of three-dimensional rotations, defined by (26),
we find that it is simple and that the metric tensor is
\[
g_{\rho \sigma}=2 \delta_{\rho \sigma}
\]
